% Chapter 4, Figure 1: model rocket in flight at time t (a) and at t + dt (b) (scan: figures/ch4/fig01.png,
% PDF 546).
% The house model rocket (tamrfig `model', units of its length 6.4, 0.6 cm each), nose up and to the right, its
% axis 60 degrees above the horizontal as on the scan, with the house exhaust plume (the 1973 spindle behind the
% tail; 0.72 of the rocket's length, 1973 about 0.6); the C.G. at 0.72 of the length from the nose (as drawn; the
% book gives no value).
% Vectors as in Chapter 1 Figures 6-8 (the house `vec'), lettered with arrows as the text sets them (the 1973
% hooked overbars; STYLE s.15, as Chapter 3 Figure 1): the velocity v vertical from the C.G.; the resultant
% external force E into the C.G. from the upper right, at 30 deg below the horizontal (1973: it ends on the body
% just below the C.G.).
% (b) is (a) a time dt later: v + dv is the resultant of v and dv (head to tail; dv 68 deg above the horizontal,
% |dv| = 0.55 |v|, as drawn); the expelled mass dm_e is the hatched slice of exhaust just behind the nozzle exit
% (hatch back, so the lines cross the 60 deg axis at 75 deg; 0.8 units deep);
% from its aft face the exhaust velocity c (relative to the rocket, backwards along the axis) and v (the same
% vector as at the C.G.) add head to tail to c + v, the slice's velocity relative to the ground (eq. (4)).
% The vectors are longer against the rocket than in 1973 (|v| = 2.9 units, 1973 about 2.0; |c| = 4.2, 1973 about
% 3.4), so the triangle at the C.G. rises clear of the nose and its label v + dv sits beside its head (the 1973
% label lies across the body); c + v then points 19 deg below the horizontal (1973 about 35 deg).
% The state lists are typeset with their two columns as lettered.
\documentclass[11pt]{standalone}
\usepackage{tamrfig}
\begin{document}
\begin{tikzpicture}[x=0.6cm, y=0.6cm]
  \def\ax{60}          % axis angle above the horizontal
  \def\L{6.4}\def\cgs{4.6}\def\plume{4.6}
  \def\vl{2.9}         % |v|
  \def\cl{4.2}         % |c|
  \def\El{4.0}\def\Ea{210}   % E: length, direction
  % \rocket: the rocket with its nose tip at the origin of the current scope
  \newcommand{\rocket}{%
    \pic[rotate=180+\ax] {rocket={model, centerline=false, plume=\plume}};
    \draw[centerline] (\ax:0.3) -- (180+\ax:\L+\plume+0.3);
    \coordinate (cg) at (180+\ax:\cgs);}
  % ---------------- (a): time t
  \begin{scope}
    \rocket
    \draw[vec] (cg) -- ++(0,\vl) node[pos=0.6, left=2pt] {$\vec{v}$};
    \draw[vec, shorten >=2.8pt] ($(cg)+(\Ea-180:\El)$) -- (cg) node[pos=0.45, below right=0pt and 0pt] {$\vec{E}$};
    \node[cg mark] at (cg) {};
    \node[anchor=north west, inner sep=0pt] at (-11,1.9) {%
      \begin{tabular}{@{}l@{\hspace{0.6em}}l@{}}
        \multicolumn{2}{@{}l}{Time $= t$}\\
        Rocket & mass $= m(t)$\\
        Rocket & velocity $= \vec{v}$\\
        Rocket & momentum $= \vec{p}$
      \end{tabular}};
    \node[panel, anchor=south east] at (3.4,-9.5) {(a)};
  \end{scope}
  % ---------------- (b): time t + dt
  \begin{scope}[shift={(0,-13.2)}]
    % the expelled mass: a hatched slice of the exhaust just behind the nozzle exit
    \def\slab{0.8}
    \begin{scope}[shift={(180+\ax:\L)}, rotate=180+\ax]
      \clip \plumeshapeb{1.3}{0.2}{\plume};
      \fill[hatch back] (0,-1) rectangle (\slab,1);   % 135 deg: 75 deg across the axis
      \draw[thin line] (\slab,-1) -- (\slab,1);
    \end{scope}
    \rocket
    \coordinate (S) at (180+\ax:\L+\slab);    % the aft face of the slice
    % at the C.G.: v, dv, v + dv
    \coordinate (vt) at ($(cg)+(0,\vl)$);
    \coordinate (dvt) at ($(vt)+(68:1.6)$);
    \draw[vec] (cg) -- (vt) node[pos=0.55, left=2pt] {$\vec{v}$};
    \draw[vec] (vt) -- (dvt) node[pos=0.5, left=2pt] {$d\vec{v}$};
    \draw[vec] (cg) -- (dvt) node[pos=1, anchor=south west, inner sep=1.5pt, xshift=1pt] {$\vec{v} + d\vec{v}$};
    \draw[vec, shorten >=2.8pt] ($(cg)+(\Ea-180:\El)$) -- (cg) node[pos=0.45, below right=0pt and 0pt] {$\vec{E}$};
    \node[cg mark] at (cg) {};
    % at the slice: c (backwards along the axis), then v, make c + v
    \coordinate (ct) at ($(S)+(180+\ax:\cl)$);
    \coordinate (cvt) at ($(ct)+(0,\vl)$);
    \draw[vec] (S) -- (ct) node[pos=0.5, below right=0pt and 0pt] {$\vec{c}$};
    \draw[vec] (ct) -- (cvt) node[pos=0.5, left=2pt] {$\vec{v}$};
    \draw[vec] (S) -- (cvt) node[pos=0.75, above left=0pt and -2pt] {$\vec{c} + \vec{v}$};
    \callout{($(180+\ax:\L+\slab/2)+(\ax-90:0.22)$)}{(-1.2,-7.25)}{$dm_e$}
    \node[anchor=north west, inner sep=0pt] at (-11,1.9) {%
      \begin{tabular}{@{}l@{\hspace{0.6em}}l@{}}
        \multicolumn{2}{@{}l}{Time $= t + dt$}\\
        Rocket & mass $= m(t) - dm_e$\\
        Rocket & velocity $= \vec{v} + d\vec{v}$\\
        Rocket & momentum $= \vec{p} + d\vec{p}$
      \end{tabular}};
    \node[panel, anchor=south east] at (3.4,-9.5) {(b)};
  \end{scope}
\end{tikzpicture}
\end{document}
